% Copyright 2026 Alan Szmyt. All rights reserved pending publication license.
\section{Discussion}
\label{sec:discussion}

The proposed framing emphasizes an inspectable chain across consented context,
semantic intent, time-indexed planning, measured acoustic realization, actual
exposure, separated response measures, and advisory updating. The bounded
review did not identify that complete governed chain in one evaluated system,
but each major technical idea has prior precedent. The contribution is
therefore a formalized combination and operational separation, not a verified
first-system claim.

Following AffectMachine-Classical, the term \emph{probability-space sculpting}
is used only as a computational metaphor \citep{agres2023affectmachine}:
semantic prompting and other controls constrain a conditional distribution over
possible generations rather than deterministically selecting one sound. The
metaphor becomes scientifically useful only if the intended controls and their
realized acoustic consequences can be operationalized and tested.

\subsection{Contribution at the design evidence level}
\label{sec:discussion-design-contribution}

Antidote's contribution at this stage is the specification of a research
object and its governance, not evidence that an adaptive audio system changes a
listener's state. Current-to-target selection, person-specific modeling,
generated affective trajectories, continuous musical feedback, interpretable
controls, and semantic planning each have prior precedent
\citep{janssen2012tune,daly2016abcmi,ehrlich2019closedloop,
agres2023affectmachine,zhang2026mindmelody}. The narrower contribution is to
formalize how consented context, semantic intent, a time-indexed plan,
adapter-specific conditioning, acoustic realization, actual exposure,
separated response constructs, provenance, and a possible later update can
remain individually inspectable. Our bounded review did not identify that
complete governed chain in one evaluated system. This supports a candidate
research-system contribution, not a global first-system claim.

Two separations carry most of the design rationale. The first is epistemic:
the record that requests a property is not the measurement of its realization,
and neither record determines what a person perceived, felt, found helpful, or
found harmful. The second concerns timing and authority: in the proposed
two-rate architecture, uncertain semantic deliberation may prepare verified
future material, while a deterministic renderer would own audio deadlines and
the immediate stop path. These separations make disagreements locatable in
principle. A plan may be accepted while an adapter visibly downgrades one
control; an artifact may satisfy its file checks while missing a requested
acoustic property; or playback may be technically continuous while the person
reports mismatch. Preserving those possibilities is more defensible than
calling a prompt, affect label, or aggregate score ``interpretable.'' It does
not guarantee that the records are understandable, correct, complete, or safe.

Three proposed components turn this separation into falsifiable design claims.
The Semantic Intent Mixer can be tested for preservation of typed controls,
exclusions, authority precedence, conflict resolution, and visible adapter
downgrades; its control values do not imply linear changes in acoustics or felt
response. The predictive horizon is a receding horizon over candidate controls,
not a forecast of the person's lived state: it can be tested by whether only the
first approved action is committed and whether verified future material meets a
frozen generation deadline. Continuity must be tested separately as
semantic-plan distance, measured acoustic-boundary distance, and waveform
scheduling; success at one layer cannot establish the others or a person's
perception of continuity. None of the three is implemented or validated end to
end.

Five evidentiary distinctions must therefore remain non-interchangeable.
Demonstrated implementation means that a named capability executes at a
reviewed source revision; at that revision, the only executable session
evidence is the rule-guided synthetic slice described in
Section~\ref{sec:design-two-rate-architecture}. Technical feasibility must
remain stage-specific: qualifying T0 records can establish only the synthetic
authority, schema, replay, interruption, recovery, hash, and provenance
invariants; T1 would additionally require real-model adherence, continuity,
latency and buffer-margin, and failure-and-recovery evidence under a frozen
model, adapter, analyzer, hardware, mutation, and seed manifest. Experiential
feasibility would require authorized H1 records showing that the procedure and
its separated response measures can be completed with acceptable burden; it
would not establish efficacy. Efficacy would require a materially different
design capable of supporting a causal outcome claim, while a mechanism claim
would require evidence identifying how an effect arose rather than merely
observing change. Neither efficacy nor mechanism is tested here.

These arguments establish only the design-level boundary described by RQ1.
Formal equations and a coherent record model cannot by themselves answer RQ2
or RQ3. The contribution would need revision if the record chain proves too
burdensome to use, if its provenance cannot expose where material mismatches
arise, if model adapters cannot preserve the approved semantic distinctions, or
if simpler records support equally accountable evaluation. Inspectability is
consequently a falsifiable design commitment, not an empirical success already
achieved.

\subsection{Moment-specific and negotiated personalization}
\label{sec:discussion-negotiated-personalization}

The person--moment unit resists treating personalization as retrieval of a
fixed profile or discovery of an objectively correct sound. A person may want
to stay with, explore, soften, or close an experience at one time and make a
different choice later. Antidote therefore represents prior evidence as
eligible context rather than destiny, a desired journey as a proposal rather
than a promised endpoint, and semantic language as person-extensible rather
than as a universal affect vocabulary. ``Moment-specific'' describes this
conditioning boundary; it does not imply that the system knows what the person
needs.

``Negotiated'' identifies an authority process rather than an optimization
result. Within the proposed framework, person-authored intent, an inspectable
plan, adapter downgrades, and any evidence-linked update proposal remain
distinct, reviewable records; edits create new revisions, and approvals or
rejections are recorded where the applicable contract permits them. Model
uncertainty is not converted into authority by being presented as
personalization. This process may make the system's assumptions easier to
contest, but additional interaction may itself create burden, demand
characteristics, or reduced trust. An accepted plan is thus evidence of a
recorded choice under the declared interface, not evidence that the plan
represents a stable preference, was freely understood, or will be helpful.

The response model must be equally conditional. Musical responses may involve
memory, expectancy, imagery, appraisal, cultural associations, and other
pathways that are not recoverable from one acoustic feature or affect label
\citep{juslin2008emotional}. Perceived expression and felt emotion can differ
\citep{zentner2008emotions}, while intensity may be welcome, unwanted, mixed,
or harmful in ways that require the person's interpretation
\citep{silverman2020complicated}. The same realized artifact may therefore
receive different accounts in different moments without one being treated as
noise merely because it conflicts with an earlier record.

The frozen H1 design can eventually examine whether those distinctions can be
collected with acceptable completeness and burden, but its P condition is a
bundle: it changes contextual input, personal semantic content, plan visibility,
editing, and approval together. A P-versus-S or P-versus-G contrast cannot
identify which component mattered, and a favorable response cannot validate an
objective need estimate. Even repeated provenance-linked records would support
only a bounded within-person account unless a later protocol prospectively
tests whether a named advisory relationship predicts new observations. Until
then, moment-conditioned personalization remains a structured hypothesis and a
human-authority requirement, not an observed advantage.

\subsection{Alternative explanations and architectures}
\label{sec:discussion-alternatives}

Any later favorable response would admit explanations other than successful
moment-specific adaptation. Generic ambient audio might suffice; a current-state
category with a fixed journey template might capture most of the difference; or
preference, familiarity, novelty, expectancy, attention, editing effort, and
interface-mediated agency might shape the report. Output quality, loudness,
setting, time of day, baseline variation, order, or carryover could also
confound the intended semantic manipulation. Repeated use may teach instrument
use or reporting rather than reveal a stable audio--response relationship. The
G, S, and P conditions preserve an initial comparison among interface bundles,
but cannot isolate components or mechanisms.

Contradictory evidence makes those alternatives substantive rather than
ceremonial. Prior systems demonstrate technically closed loops and
person-specific or affective control, yet their evaluations are small,
heterogeneous, or directed at different constructs. A recent preprint of a
preliminary brain--computer musical-interface evaluation reported nonsignificant
target-emotion and time effects, with person and trial variation dominating its
observed variance \citep{monroy2026minimalist}. A systematic review likewise
found heterogeneous
music-neurofeedback methods and no consensus success measure
\citep{sayal2025musicloop}; a scoping review of commercial music-based digital
therapeutics found insufficient direct evidence for real-world effectiveness
or superiority over generic music \citep{venkatesan2026mdt}. These sources do
not show that an individual effect is impossible. They do show why loop closure,
market presence, and a beneficial response cannot substitute for one another.

The proposed architecture is also only one arrangement of the technical
problem, and four simpler baselines deserve direct comparison. A population
affect-label pipeline could map a small shared state vocabulary to fixed audio
rules without personal context. A static recommender could select or sequence
existing recordings, as prior current-to-target systems do
\citep{janssen2012tune,mallik2022effects}. Direct physiological conditioning
could drive continuous controls without a negotiated semantic layer, following
earlier brain--computer music systems \citep{daly2016abcmi,
ehrlich2019closedloop}. A one-shot generator could render the complete journey
offline for human review, avoiding a future-audio buffer and mid-session model
deadline altogether. Interpretable rule-based and probabilistic generation
\citep{agres2023affectmachine} and human-authored segments with a fixed
transition engine provide further alternatives. These arrangements may reduce
latency, unsupported controls, interaction burden, or transition uncertainty
at the cost of a smaller or less inspectable generation space. Antidote's
two-rate proposal should be preferred only if explicit deadlines,
verification, and fallbacks measurably improve the declared technical
properties without obscuring person authority.

Several outcomes would count against the proposed arrangement: adapters that
frequently downgrade accepted intent, a buffer margin that cannot be maintained
on the frozen hardware, fallbacks that dominate listening, semantic or acoustic
boundary measures that do not correspond to reviewable failures, or an
authority path that does not halt reliably. Even technical success would leave
the response question open. Small boundary distances, accurate tempo, or an
uninterrupted file cannot establish perceived continuity, fit, usefulness,
safety, or benefit. Alternative architectures should therefore be compared
under the same provenance, exposure, missingness, burden, and non-positive
reporting rules rather than by selecting whichever system produces the most
persuasive demonstration.

\subsection{Bounded future research program}
\label{sec:discussion-future-program}

The next evidence should be acquired within the stages and gates already frozen
by the protocol rather than by expanding the interpretation of the present
record. At D0, the immediate question is whether every RQ1 construct and
governed transformation or authority boundary resolves to a versioned record,
equation, contract, or explicit unavailable state; satisfying that traceability
check would not constitute an independent implementation audit. At T0, a
qualifying synthetic package would have to preserve complete denominators and
demonstrate the declared authority, schema, replay, interruption, recovery,
hash, and provenance invariants. Only an explicit claim-ledger promotion could
turn that package into a synthetic technical result, and no T0 result could
describe a human exposure or response.

T1 is the next test of the proposed generation path. Before execution it would
freeze the model, adapter, hardware profile, analyzer set, mutation manifest,
seed policy, output-rights review, hardware- and adapter-specific thresholds,
and run plan.
Its 39-attempt matrix, when every declared control is supported, would test
one-factor adherence, adjacent-segment continuity, latency and buffer margins,
and injected failure and recovery. A technically feasible claim would require
the reserved source package, complete terminal-state accounting, reproducible
analysis, and review against those frozen thresholds. Any resulting claim
would remain bounded to that manifest. Unsupported controls would remain
explicitly unsupported and reduce the applicable attempt denominator; missed
deadlines against frozen thresholds, integrity failures, or unreliable stopping
would require the corresponding technical claim to be downgraded or rejected,
or the design to be revised. Passing T1 would remain specific to the tested
model and hardware and would not establish perceived continuity or a beneficial
response.

Optional real-time sensing and genuinely continuous generation are separate
later studies, not capabilities assumed by H1. Every sensor stream would need
its own consent, calibration, missingness, uncertainty, synchronization, and
person-correction contract; its estimate could tailor a proposal but could not
define the person's state. A streaming generator would likewise need a new
qualified adapter and deadline study demonstrating that variable-latency model
work remains outside the audio callback and that stop does not depend on model
completion. Failure at either gate should default to manual input and, only
after its own implementation and qualification, generate-ahead playback; it
must not silently strengthen inference or timing claims.

H1 remains \texttt{blocked-no-collection-authority}: all eight activation
gates---real-model/T1 qualification, privacy, consent, independent review,
assignment, instrument, analysis, and pilot stop/go---remain unsatisfied. If
separately authorized, its six-block G/S/P series targets 18 completed exposures
within an absolute ceiling of 24 scheduled attempts and can test whether the
governed records are collectable, correctable, provenance-linked, and usable
with the prespecified completeness, burden, and safety-event accounting. It can
describe exploratory within-person bundle contrasts while preserving missing,
null, opposite-direction, interrupted, burdensome, harmful, and adverse
observations. It cannot isolate personalization components, estimate a
population effect, establish safety, or answer whether advisory learning is
useful. Its progression rule decides only whether another bounded study may be
acceptable; ``proceed'' is not an outcome claim.

Testing RQ3 would consequently require a new, prospectively frozen protocol
after H1, not reuse of the H1 observations both to discover and confirm a
preferred relationship. Such a protocol would need to name the eligible semantic,
acoustic, exposure, and response variables; specify the response model,
uncertainty, forgetting and contradiction rules; separate model development
from evaluation on later observations; and retain human review of every
proposed update. A useful advisory claim would require prospective evidence
that a frozen proposal rule adds information for that person's new records
relative to declared non-adaptive alternatives, while retaining burden,
mismatch, harm, and missingness. Failure to improve prospective prediction, or
improvement purchased through unacceptable burden or risk, would count against
the claim.

Only after those within-person gates should multi-participant work be
considered. A later study would require a separate ethics
determination, participant and setting rationale, sample-size justification,
device and model qualification, accessible procedures, prespecified comparison
conditions, and an analysis that distinguishes person variation from moment,
order, and trial variation. It would also need to test rather than assume that
the semantic vocabulary, control mappings, and response measures transfer
across people and contexts. No stage described here is a treatment or
neurological-mechanism evaluation. A claim may be promoted only when its named
source package and review gate support that promotion; otherwise the relevant
slot remains unavailable or the candidate remains unchanged, is downgraded, or
is rejected. The unresolved validity, emotional and hearing safety, privacy,
rights, accessibility, and independent-review gates in
Section~\ref{sec:limitations} constrain every future direction described here;
publication of the design cannot satisfy them.
